\documentclass[10pt,a4paper]{amsart}
\usepackage[margin=19mm]{geometry}
\usepackage{amsmath,amssymb,mathtools}
\usepackage[scr=rsfs,cal=euler]{mathalfa}
\usepackage{xfrac}
\usepackage[unicode,hidelinks,bookmarks=false]{hyperref}
\newcommand{\ie}{i.e.\ }
\newcommand{\cf}{cf.\ }
\newcommand{\resp}{respectively\ }
\newcommand{\Subsection}{\subsection}
\providecommand{\nobreakdash}{\nobreak\hbox{-}}
\setlength{\emergencystretch}{2em}
\multlinegap0pt
\usepackage{bm}
\usepackage{bbm}
\usepackage{dsfont}
\usepackage{xspace}
\usepackage{cancel}
\usepackage{mathtools}
\usepackage{amsthm}
\newtheorem{proposition}{Proposition}
\newcommand{\vect}[1]{\ensuremath{\bm{#1}}}
\title[A Conjugate Gradient Formulation of the EnKF Algorithm]{A Conjugate Gradient Formulation of the EnKF Algorithm (Proposition 3)}
\author{Sanghyun Lee, Zhengqi Liu, Jonathan Valyou, Ludmil Zikatanov}
\date{Source: arXiv:2606.19224v2 (CC BY 4.0)}
\begin{document}
\maketitle

\noindent\emph{Source and licence.} A Conjugate Gradient Formulation of the EnKF Algorithm. Version arXiv:2606.19224v2 (CC BY 4.0), distributed under the Creative Commons Attribution 4.0 International licence, as recorded in the arXiv licence metadata for this exact version and shown on the abstract page https://arxiv.org/abs/2606.19224v2. Full rights notice: licenses/arxiv-2606.19224-CC-BY-4.0.txt.\par\smallskip

\noindent\textit{Editorial scope.} Counted unit: the Proposition printed as Proposition 3 of the article (source0.tex lines 671--680, source label \texttt{prop:3}), delivered together with the article's own complete original proof of it (source0.tex lines 683--787, one \texttt{proof} environment of 105 lines). No mathematics is written, repaired or re-derived: statement and proof are the article's own text line for line. Required context retained uncounted: none. Every definition, display and label of those lines is retained unchanged. Omitted from this package and not redistributed: 1--670, 681--682, 788--1747. Nothing inside a retained excerpt is omitted or altered: every retained body is delivered exactly as the article's lines are written, so no omission receipt of any kind accompanies this package. Citations: the retained excerpts carry no citation command at all, so no bibliography entry of the article is transcribed and no reference is redistributed. Reference adaptation: none was needed and none was made; every reference inside the delivered unit resolves inside it, so no unresolved reference or citation is produced. Printed slips kept literal: no printed word, symbol, hypothesis or number of the counted statement or of its proof was altered. Layout and class: the article's own class is \texttt{elsarticle} with packages including amsmath, amsfonts, verbatim, amsbsy, multirow, color, graphicx, amsthm, makecell, tikz, relsize, scalefnt, tabularx, amssymb; the wrapper is the lane's pinned \texttt{amsart} editorial wrapper at 10pt on A4, which restates the theorem-like environments the retained text needs and adds the article's own macro definitions that the retained text needs (\texttt{\textbackslash vect}), with extra wrapper packages (bm, bbm, dsfont, xspace, cancel, mathtools). The delivered numbering is the unit's own. Assets: no retained span contains a figure environment, a graphics-inclusion command, a PDF-page inclusion, a table or a TikZ picture; no image file is copied into this package, the package declares no asset dependency and no image file is redistributed with it. Licence: version arXiv:2606.19224v2 of this article is distributed under the Creative Commons Attribution 4.0 International licence, as recorded in the arXiv licence metadata for this exact version and shown on the abstract page https://arxiv.org/abs/2606.19224v2. A full rights notice is delivered at licenses/arxiv-2606.19224-CC-BY-4.0.txt. Characters: every retained excerpt is pure ASCII, so the delivered engine, XeLaTeX with its default Latin Modern text font, reports no missing character.
\begin{proposition}
Let $1\leq r<p$.  Let the ordered eigenvalues of 
    $\Psi_{z}$ be $\lambda_1 \geq \lambda_2 \geq ...\geq \lambda_{b-1} \geq \lambda_b \geq \lambda_{b+1} \geq  ... \geq \lambda_{a-1} \geq \lambda_a \geq \lambda_{a+1} \geq  ... \geq \lambda_{p-1} \geq \lambda_p$. Let there be $l$ unique eigenvalues of $\Psi_{z}$ that are less than $\lambda_a$ or greater than $\lambda_b$ denoted $\mu_i$ for $i=1,...,l$ in a set $\sigma_0(\Psi_{z})$.  Additionally, let all other unique eigenvalues of $\Psi_{z}$ ranging from $\lambda_a$ to $\lambda_b$ be in a set $\sigma_1(\Psi_{z})$.  Then, if $r \geq l$, there exists the following error bound on a rank $r$ CGD approximation of $\Psi_{z}^{-1}$:
    \begin{align}
    ||\Psi_{z}^{-1} - \Sigma_r D_r^{-1} \Sigma_r^T||_{\Psi_{z}} \leq 2C(\frac{\sqrt{\operatorname{cond}(\Psi_{z})_{eff}} - 1}{\sqrt{\operatorname{cond}(\Psi_{z})_{eff}} +1})^{r-l} \| \vect{e}_0\|_{\Psi_{z}}, \\
    C = \max\limits_{\substack{\lambda \in \sigma_1(\Psi_{z})}} \prod_{i=1}^{l} | 1-\frac{\lambda}{\mu_i}|
    \end{align}
where $\operatorname{cond}(\Psi_{z})_{eff}$ is the effective condition number of $\Psi_{z}$ defined as $\frac{\lambda_b}{\lambda_a}$, $\Sigma_r \in \mathbb{R}^{p \times r}$ contains the first $r$ eigenvectors of $\Psi_{z}$, $D_r^{-1}$ is a diagonal matrix with the first $r$ eigenvalues along the diagonal, and $\vect{e}_0$ is the initial CGD approximation error.
\label{prop:3}
\end{proposition}

\begin{proof}
Since $\Psi_{z}$ is symmetric positive definite, it is diagonalizable into a matrix $D_p$ by its eigenmatrix.  Therefore,
\begin{equation}
D_p = \begin{bmatrix}
 D_r & 0\\
 0 & D_q
\end{bmatrix} = \begin{bmatrix}
 \Sigma_r^T \Psi_{z} \Sigma_r & 0\\
 0 & \Sigma_q^T \Psi_{z}^T \Sigma_q
\end{bmatrix} 
= \Sigma_p^T \Psi_{z} \Sigma_p,
\label{eq:diagonalization}
\end{equation}
where we define $\Sigma_p \in \mathbb{R}^{p \times p}$ as all eigenvectors of $\Psi_{z}$ in the columns, $\Sigma_r \in \mathbb{R}^{p \times r}$ as the first $r$ eigenvectors of $\Psi_{z}$ in the columns, and $\Sigma_q \in \mathbb{R}^{p \times (p-r)}$ as the last $p-r$ eigenvectors of $\Psi_{z}$ in the columns.

Since $D_p$ is diagonal with non-zero entries, $D_p$ is invertible, so we can apply the inverse to both sides of \eqref{eq:diagonalization}:
\begin{equation}
    \Psi_{z}^{-1} = \Sigma_p D_p^{-1} \Sigma_p^T = \Sigma_r D_r^{-1} \Sigma_r^T + \Sigma_q D_q^{-1} \Sigma_q^T.
\end{equation}
Therefore, we can write an $r$-rank approximation of $\Psi_{z}^{-1}$ as
\begin{equation}
(\Psi_{z}^{-1})_r = \Sigma_r D_r^{-1} \Sigma_r^T,
\end{equation}
and the corresponding error of the $r$-rank approximation solution to $\Psi_{z}K^T = \Psi_{x}^T$ is
\begin{equation}
    \vect{e}_r = K^T - (K^T)_r = \Psi_{z}^{-1}\Psi_{x}^T - (\Psi_{z}^{-1})_r\Psi_{x}^T
    \label{eq:e_r_def}
\end{equation}


Since we are utilizing CGD to make the $r$-rank approximation, we can rather rewrite the approximation error as 
\begin{equation}
    \vect{e}_r = \vect{e}_0 + \text{span}\{\Psi_{z}\vect{e}_0, (\Psi_{z})^2\vect{e}_0,...,(\Psi_{z})^r\vect{e}_0\}
    = (I + \sum_{j=1}^{r} \omega_j(\Psi_{z})^j)\vect{e}_0,
    \label{e_r_def_rewrite}
\end{equation}
where $\omega_j$ are the weights to cover the generated Krylov Space and $\vect{e}_0$ is the initial CGD approximation error.  
We can rewrite \eqref{e_r_def_rewrite} in terms of a matrix polynomial expression:
\begin{equation}
\vect{e}_r = \mathcal{P}_r(\Psi_{z})\vect{e}_0,
\label{e_r_mat_poly}
\end{equation}
where $\mathcal{P}_r(\Psi_{z}) = I + \sum_{j=1}^{r} \omega_j(\Psi_{z})^j$.


With the given eigenvalue $\lambda$ and its corresponding eigenvector $\vect{v}$ we obtain,
\begin{equation}
\mathcal{P}_r(\Psi_{z})\vect{v}=\mathcal{P}_r(\lambda)\vect{v},
\end{equation}
and we can rewrite \eqref{e_r_mat_poly} as
\begin{equation}
\vect{e}_r= \sum_{j=1}^{r}\mathcal{P}_r(\lambda_j)\vect{e}_0.
\end{equation}

Applying the definition of the energy norm of $\Psi_{z}$, we get
\begin{equation}
    \| \vect{e}_r\|_{\Psi_{z}}^2 = \sum_{j=1}^{r} [\mathcal{P}_r(\lambda_j)]^2 \| \vect{e}_0\|_{\Psi_{z}}^2.
\end{equation}


To begin constructing an upper bound, we want to choose the eigenvalue that maximizes the right side.  This is dependent upon the chosen polynomials $\mathcal{P}_r$:
\begin{equation}
    \| \vect{e}_r\|_{\Psi_{z}}^2 \leq \max\limits_{\substack{j=1,...,r}} [\mathcal{P}_r(\lambda_j)]^2 \| \vect{e}_0\|_{\Psi_{z}}^2.
    \label{eq:constr_upper_bound}
\end{equation}

    For the eigenvalues $\mu_i$ for $i=1,...,l$ in the set $\sigma_0(\Psi_{z})$, a matrix polynomial $\mathcal{R}_l(\lambda)$ can be defined.  We specifically choose a matrix polynomial that minimizes the effect of extremal eigenvalues contained in $\sigma_0(\Psi_{z})$:

    \begin{equation}
    \mathcal{R}_l(\lambda) = \prod_{i=1}^{l} | 1-\frac{\lambda}{\mu_i}|.
    \label{eq: R_l}
    \end{equation}

    For the eigenvalues $[\lambda_a,\lambda_b]$ in the set $\sigma_1(\Psi_{z})$, a matrix polynomial $\mathcal{Q}_{r-l}(\lambda)$ can be defined.  It is assumed that these eigenvalues are evenly spread out.

    Since these two sets distinctly contain all of the eigenvalues of $\Psi_{z}$, then the matrix polynomials have the following relationship:

    \begin{equation}
    \mathcal{P}_r(\lambda) = \mathcal{R}_l(\lambda)\mathcal{Q}_{r-l}(\lambda).
    \end{equation}

    Substituting into \eqref{eq:constr_upper_bound} and finding the matrix polynomial that minimizes the right side:
    \begin{align}
    \| \vect{e}_r\|_{\Psi_{z}}^2 
    &\leq \min\limits_{\substack{\mathcal{R}_l,\mathcal{Q}_{r-l}}} \max\limits_{\substack{\lambda}}[\mathcal{R}_l(\lambda)]^2[\mathcal{Q}_{r-l}(\lambda)]^2 \| \vect{e}_0\|_{\Psi_{z}}^2 \\
    &\leq \min\limits_{\substack{\mathcal{R}_l}} \max\limits_{\substack{\lambda}}[\mathcal{R}_l(\lambda)]^2\min\limits_{\substack{\mathcal{Q}_{r-l}}} \max\limits_{\substack{\lambda \in \sigma_1(\Psi_{z})}}[\mathcal{Q}_{r-l}(\lambda)]^2 \| \vect{e}_0\|_{\Psi_{z}}^2.
    \end{align}

    Substituting in the chosen matrix polynomial $\mathcal{R}_l(\lambda)$ from \eqref{eq: R_l} and choosing a Chebyshev polynomial form for $\mathcal{Q}_{r-l}(\lambda)$ due to the equal spacing assumption, we obtain

    \begin{equation}
    \| \vect{e}_r\|_{\Psi_{z}}^2 \leq [\max\limits_{\substack{\lambda \in \sigma_1(\Psi_{z})}} \prod_{i=1}^{l} | 1-\frac{\lambda}{\mu_i}|]^2[(2\frac{\sqrt{\frac{\lambda_b}{\lambda_a}} - 1}{\sqrt{\frac{\lambda_b}{\lambda_a}} +1})^{r-l}]]^2 \| \vect{e}_0\|_{\Psi_{z}}^2 \text{, and}\\
    \end{equation}

    \begin{equation}
    \| \vect{e}_r\|_{\Psi_{z}}^2 \leq [C]^2
    [(2\frac{\sqrt{\operatorname{cond}(\Psi_{z})_{eff}} - 1}{\sqrt{\operatorname{cond}(\Psi_{z})_{eff}} +1})^{r-l}]^2 \| \vect{e}_0\|_{\Psi_{z}}^2 .\\
    \end{equation}

    Taking the square root of both sides gives an upper limit error bound with respect to the energy norm:

    \begin{equation}
    ||\vect{e}_r||_{\Psi_{z}} \leq 2C(\frac{\sqrt{\operatorname{cond}(\Psi_{z})_{eff}} - 1}{\sqrt{\operatorname{cond}(\Psi_{z})_{eff}} +1})^{r-l} \| \vect{e}_0\|_{\Psi_{z}}.
    \end{equation}
\end{proof}

\end{document}
